% SOURCE_AUTHORITY: sga5_fr_workpass.tex, lines 6861--6876 (LF-normalized unit).
% SOURCE_SHA256: 791F4EFFC5E02832D5D77ED03518C8156D6F07E4C8238B03545DB93D883FBB28
\subsubsection*{2.3.3}
Supongamos que $S$ es un subconjunto de $\Fl(\cE)$ que permite un cálculo de fracciones por la derecha. Se deduce entonces de la parte (d) de 2.3.1 que, para todo objeto $X$ de $\cE$, la categoría $(S/X)^\circ$, opuesta a la categoría de las flechas sobre $X$ que pertenecen a $S$, es filtrante (SGA 4 I 2.7). Consideremos entonces la categoría $\cE_S$ descrita del modo siguiente:
\begin{itemize}
\item los objetos de $\cE_S$ son los objetos de $\cE$;
\item si $X$ e $Y$ son dos objetos de $\cE$,
\[
\Hom_{\cE_S}(X,Y)=\varinjlim_{(S/X)^\circ}\Hom_{\cE}(\source(\cdot),Y).
\]
\end{itemize}
Si $P_S$ designa el funtor evidente de $\cE$ en $\cE_S$, el par $(P_S,\cE_S)$ es solución del problema universal descrito al comienzo del apartado. En este caso convendremos en denotar por $\cE(S^{-1})$ esta solución particular.

\subsubsection*{2.3.4}
Sea ahora $S$ un subconjunto de $\Fl(\cE)$ que permite un cálculo de fracciones por la izquierda. Entonces el funtor $P_S:\cE\to\cE(S^{-1})$ conmuta con los límites inductivos finitos. En particular, si $\cE$ es aditiva, $\cE(S^{-1})$ es aditiva y $P_S$ es un funtor aditivo.

\subsubsection*{2.3.5}
Si la categoría $\cE$ es abeliana y $S$ es un subconjunto de $\Fl(\cE)$ que permite a la vez un cálculo de fracciones por la derecha y un cálculo de fracciones por la izquierda, la categoría $\cE(S^{-1})$ es abeliana y el funtor $P_S$ es exacto. Además, el par $(P_S,\cE(S^{-1}))$ es un cociente de $\cE$ por la subcategoría abeliana gruesa $P_S^{-1}(0)$.
